\chapter{Criterio de política editorial sobria de publicación}

\section{Política editorial sobria}
\begin{definition}
Se llama política editorial sobria de publicación al conjunto disciplinado de reglas mediante las cuales el sistema MADELEINE decide qué piezas publica, bajo qué formato, con qué grado de amplitud y con qué reservas explícitas.
\end{definition}

\begin{observation}
La política editorial sobria no equivale a censura ni a apertura total. Su función es impedir tanto la inflación prematura del corpus como su compresión abusiva en salidas externas no defendibles.
\end{observation}

\section{Condición mínima de política editorial}
\begin{definition}
Existe política editorial sobria cuando pueden declararse simultáneamente:
\begin{enumerate}
  \item criterios legibles de selección entre material interno, observado y público;
  \item reglas mínimas de formato para ficha, dossier, mapa o publicación integral;
  \item condiciones explícitas de reserva, observación o no salida;
  \item continuidad visible con la genealogía metodológica del corpus;
  \item una proporción defendible entre sobriedad, claridad y amplitud externa.
\end{enumerate}
\end{definition}

\begin{observation}
No se exige todavía una clausura definitiva de toda fase futura del sistema. Basta con una política suficiente para evitar arbitrariedad editorial inmediata.
\end{observation}

\section{Política y jerarquía del corpus}
\begin{definition}
Toda política editorial sobria debe preservar la diferencia funcional entre:
\begin{enumerate}
  \item núcleo fundacional;
  \item canon intermedio;
  \item capa metodológica;
  \item claims públicos;
  \item material observado;
  \item material reservado.
\end{enumerate}
\end{definition}

\begin{observation}
Esta cláusula impide que la publicación trate como equivalentes piezas que pertenecen a niveles distintos de madurez o función.
\end{observation}

\section{Política y formatos externos}
\begin{definition}
La política editorial sobria puede articular la salida externa mediante:
\begin{enumerate}
  \item fichas públicas mínimas;
  \item dossiers externos de claims seleccionados;
  \item mapa público mínimo del corpus;
  \item publicación integral sobria del sistema;
  \item anexos metodológicos o reservas explícitas cuando resulten necesarios.
\end{enumerate}
\end{definition}

\begin{observation}
La política no debe favorecer un formato por costumbre o atractivo retórico, sino por ajuste real entre forma externa y madurez metodológica de la pieza.
\end{observation}

\section{Condición mínima de no política}
\begin{definition}
No debe declararse política editorial sobria cuando la salida externa dependa de intuiciones no explicitadas, criterios variables sin justificación o simplificaciones que borren la jerarquía interna del corpus.
\end{definition}

\begin{observation}
La ausencia de política editorial explícita no invalida el sistema, pero sí debilita su proyección pública y su coherencia de publicación.
\end{observation}

\section{Cláusula de alcance}
\begin{observation}
El criterio de política editorial sobria de publicación fija la primera regla explícita de salida externa del sistema MADELEINE. No constituye todavía una clausura exhaustiva de todas las decisiones editoriales futuras del corpus.
\end{observation}

